% 出席確認クイズ 第05回　業務プロセスとAIによる自動化
%   attendance/attendance05.html から生成（解説は本ガイド用に加筆）。

\quizq{1}{RPAの説明として適切なものはどれでしょうか?}%
  {ネットワーク機器の一種}%
  {工場で働く物理的な産業用ロボット}%
  {\correct{定型的なPC操作をソフトウェアロボットで自動化するしくみ}}%
  {会計基準の名称}%
  {正解は (c)。RPA（Robotic Process Automation）は、パソコン上の定型作業（転記、集計、入力）をソフトウェアロボットで自動化する技術です。(b) の物理的な産業用ロボットとは別物です。ルールが明確で手順が変わらない作業に向きます。}

\quizq{2}{RPAと生成AIの違いとして適切なものはどれでしょうか?}%
  {RPAは自然な文章を生成できる}%
  {\correct{RPAは決まった手順の反復に強く、生成AIは非定型の文章生成や判断支援に向く}}%
  {生成AIは定型処理しかできない}%
  {両者はまったく同じもの}%
  {正解は (b)。RPA はルール化された定型作業の反復に強く、生成AIは文章作成・要約・判断支援など非定型の作業に向いています。両者は競合ではなく、RPA が生成AIの出力を基幹システムに入力するといった組み合わせで使われます。(a)(c)(d) は特徴を取り違えています。}

\quizq{3}{「AIエージェント」の説明として適切なものはどれでしょうか?}%
  {人間の営業担当者のこと}%
  {一つの質問に一度だけ答えるしくみ}%
  {AIの利用契約書のこと}%
  {\correct{目標に向けて自律的に計画し、ツールを使って複数の手順を実行するAI}}%
  {正解は (d)。AIエージェントは、目標を与えられると自ら計画を立て、検索やシステム操作などのツールを使って複数の手順を実行し、結果を確認しながら進めるAIです。(b) はチャットの一問一答で、エージェントではありません。自律性が高い分、第13回で扱うガバナンスがより重要になります。}

\quizq{4}{BPM(ビジネスプロセス管理)で最初に行うこととして適切なものはどれでしょうか?}%
  {担当者を減らす}%
  {すべての業務を即座に自動化する}%
  {\correct{業務の流れを可視化する}}%
  {新しいシステムを購入する}%
  {正解は (c)。BPM（ビジネスプロセス管理）は、現状（As-Is）の可視化から始めて、分析・改善（To-Be）・実行・監視を繰り返します。可視化をせずに自動化すると、無駄な手順を高速化するだけになります。(a)(b)(d) は手段を目的と取り違えた例です。}

\quizq{5}{自動化の対象を選ぶ際の考え方として適切なものはどれでしょうか?}%
  {すべての業務を無条件に自動化する}%
  {\correct{判断が曖昧で例外の多い業務は人が担い、定型で量の多い業務から自動化する}}%
  {例外の多い業務から優先して自動化する}%
  {自動化は一切行わない}%
  {正解は (b)。自動化の効果は「定型性 × 量」で決まります。定型で件数の多い業務から自動化し、判断が曖昧で例外の多い業務は人が担うのが基本です。(c) は優先順位が逆で、(a)(d) は極端です。}
